\chapter{Conclusion et perspectives}

\section{Bilan du travail}
Ce memoire a presente \logminer, un prototype multi-agents pour la detection d'anomalies dans des journaux systemes et reseaux heterogenes. Le travail a couvert l'analyse des formats de logs, la conception d'une architecture modulaire, l'implementation d'un pipeline complet, la mise en place d'un dashboard et l'evaluation sur plusieurs familles de donnees.

Le prototype valide une chaine complete allant de la collecte a la visualisation. Les journaux sont parses, normalises, routes vers une famille, traites par un modele compatible, puis transformes en anomalies candidates et incidents. Cette approche evite de reduire le probleme a un simple score de detection.

\section{Verification des objectifs}
Les objectifs definis dans le document directeur sont couverts. Les types de journaux ont ete identifies et categorises. Les techniques classiques et modernes ont ete analysees. Une architecture d'agents specialises a ete proposee et implementee. Des mecanismes d'apprentissage supervises et non supervises ont ete integres. Un dashboard a ete realise. Le systeme a ete teste sur des journaux locaux, des exports Wazuh et plusieurs datasets publics. Enfin, les performances ont ete evaluees par F1-score, latence, ressources et comparaison operationnelle prudente.

\section{Apports}
Les apports principaux sont :
\begin{itemize}
  \item une architecture multi-agents claire et extensible ;
  \item une normalisation commune pour des journaux heterogenes ;
  \item un routage multi-modeles auditable ;
  \item une evaluation multi-source combinant logs systemes, SIEM, reseau et scenarios de robustesse ;
  \item un dashboard permettant l'analyse, les filtres, les decisions et l'audit ;
  \item un pack de reproduction comprenant figures, tableaux, captures et CSV.
\end{itemize}

\section{Limites}
Plusieurs limites doivent etre reconnues. D'abord, la distribution validee est principalement logique et locale. Redis Streams et MQTT montrent une trajectoire d'evolution, mais ne constituent pas encore une preuve de scalabilite industrielle. Ensuite, les anomalies non supervisees restent des signaux candidats. Elles doivent etre interpretees par correlation, contexte et validation humaine.

Les scores tres eleves obtenus sur certains datasets supervises doivent aussi etre lus avec prudence. Ils valident le comportement du modele dans le cadre experimental, mais ne prouvent pas une generalisation automatique a toutes les infrastructures. Enfin, la reduction des faux positifs et l'apprentissage continu restent des axes ouverts.

\section{Perspectives}
Les perspectives principales sont les suivantes :
\begin{enumerate}
  \item deployer le systeme sur plusieurs machines avec workers Redis et files persistantes ;
  \item ajouter des tests de charge prolonges et des scenarios de panne volontaire d'agents ;
  \item integrer un apprentissage actif a partir des decisions analyste ;
  \item enrichir le parsing par des templates de logs de type Drain ;
  \item reduire les faux positifs par calibration des seuils et contexte temporel ;
  \item relier les incidents a des tactiques MITRE ATT\&CK ;
  \item produire une integration plus stricte avec un SIEM existant.
\end{enumerate}

\section{Conclusion generale}
\logminer{} montre qu'une architecture locale, modulaire et interpretable peut soutenir la detection d'anomalies dans des journaux heterogenes. Sa force principale n'est pas de remplacer les outils existants, mais d'organiser une chaine coherentement auditable entre logs bruts, modeles specialises, anomalies candidates, incidents et decisions analyste.

Le prototype constitue ainsi une base solide pour un systeme de cyber-audit leger adapte aux environnements disposant de ressources limitees. Les resultats obtenus justifient une poursuite vers la scalabilite, l'integration SOC et l'apprentissage continu, tout en conservant une interpretation prudente des alertes produites.
